\section{Problématique}
\label{sec:problematique}

\subsection{Énoncé du problème}

Les équipes de marketing de performance opèrent dans un environnement à
\textbf{forte intensité décisionnelle}~: chaque jour, des dizaines de
micro-décisions doivent être prises (quel angle tester, quel créatif mettre
en pause, quel budget réallouer). Or, ces décisions sont aujourd'hui prises
sur la base d'informations \emph{fragmentées}, \emph{retardées} et
\emph{partiellement perdues}.

Plus précisément, on identifie quatre douleurs récurrentes~:

\begin{enumerate}
    \item \textbf{Fragmentation des données.} Les informations relatives
    à un même produit sont éparpillées dans 4 à 6 outils distincts. Un
    marketeur passe en moyenne 30~\% de son temps à \emph{retrouver}
    l'information avant de pouvoir l'\emph{exploiter}.

    \item \textbf{Calcul manuel du ROI.} Le retour sur investissement est
    calculé hebdomadairement dans des feuilles Excel, ce qui interdit toute
    réaction rapide à une campagne déficitaire.

    \item \textbf{Perte de la connaissance créative.} Les angles
    psychologiques (Problème/Solution, Bénéfice, Offre Irrésistible) et les
    \emph{hooks} qui ont historiquement bien performé ne sont pas
    capitalisés~: chaque nouveau produit repart d'une page blanche.

    \item \textbf{Absence d'aide à la décision en temps réel.} Aucun outil
    grand public ne propose à la fois une vue consolidée des données
    marketing \emph{et} un assistant conversationnel capable de répondre à
    des questions naturelles sur ces données.
\end{enumerate}

\subsection{En quoi le plugin DeepSeek résout-il ce problème~?}

Le chatbot intégré n'est pas un simple gadget conversationnel. Avant
d'envoyer la question de l'utilisateur à l'API DeepSeek, le backend Django
\textbf{enrichit le prompt système} avec un instantané des données les plus
pertinentes du CRM~: top campagnes des 30~derniers jours, ROI courant,
angle le plus performant, alertes budgétaires. L'assistant répond donc avec
des \emph{faits}, pas avec des généralités.

Concrètement, là où un marketeur devait auparavant ouvrir trois onglets et
faire un calcul mental pour répondre à \emph{«~Quel angle dois-je
prioriser cette semaine~?~»}, il obtient désormais une réponse argumentée en
moins de cinq secondes, dans la langue de son choix (EN/FR/AR). C'est cette
capacité à \textbf{transformer des données structurées en conseil
actionnable} qui constitue la valeur ajoutée principale du chatbot.

\subsection{Périmètre fonctionnel ciblé}

Le tableau~\ref{tab:perimetre} synthétise les fonctionnalités
\emph{incluses} et \emph{exclues} du périmètre du projet académique.

\begin{table}[H]
    \centering
    \begin{tabularx}{\textwidth}{lXX}
        \toprule
        & \textbf{Inclus} & \textbf{Exclu (perspectives futures)}\\
        \midrule
        Pipeline & Kanban Planning / Active / Évalué & Synchronisation
        Trello / Notion\\
        Assets & Angles, Créatifs (upload), Hooks & Génération automatique
        de visuels (Stable Diffusion)\\
        Finance & Saisie manuelle dépenses / revenus, ROI & Connexion
        directe Meta / Google Ads API\\
        IA & Chatbot DeepSeek contextualisé & Agents autonomes,
        function-calling avancé\\
        ML & Catégorisation Faible/Moyen/Élevé & Prévision time-series du
        budget\\
        Auth & Auth Django par défaut & SSO Google / OAuth Meta\\
        \bottomrule
    \end{tabularx}
    \caption{Périmètre fonctionnel du projet}
    \label{tab:perimetre}
\end{table}
